\section{Conclusion}
\label{sec:conclusion}

The identity-holonomy signal survives exact scrutiny, but its correct home is
narrower than initially hoped.  Chronological replay proves that the two
length-six cycles and the length-twenty-four relation have raw product $I_4$.
Exact integer cone certificates simultaneously prove that none of their
cyclic phases lies in the tested positive Rauzy length or transfer domains.
Positive raw covariant witnesses for the two shorter words show why the
separation of $B^{-\mathsf T}$, reversed $B^{\mathsf T}$, and $B$ is
essential.

At the conceptual level, categorical inversion forces an additive roof to
change sign.  Matrix identity does not turn a kernel word into a unit arrow,
and it kills the projective normalizer only on a domain where the complete
chronological path is defined.  If such a domain is artificially enlarged,
the return map is the identity and its fixed set is fully neutral.  If
instead one assigns positive symmetric edge length, one obtains a valid and
useful non-backtracking graph suspension---but a different dynamical system.

The operator analysis gives the same boundary in another language.  Bounded
faithful inverse branches on one infinite-dimensional space cannot be compact
or nuclear blocks.  Anisotropic, microlocal, and pinning theories can escape
that theorem only by supplying genuine hyperbolicity and distinct
polarizations.  The finite group-von-Neumann trace does not need ordinary
compactness and therefore retains the exact C29 determinant coefficients.

Accordingly, the tested AGY promotion is rejected at the first two
\RouteA{} gates.  This is a useful negative result: it closes an entire
formal-inverse route without a longer word or prime scan, while preserving
the finite algebraic object.  The next credible construction should place the
identity-holonomy cocycle in a fibre over an independently certified
hyperbolic base and prove an all-word determinant-tail theorem.  That move
changes the dynamical architecture rather than relabelling a failed positive
cone.
